\chapter{Substitusi Nukleofilik}\label{ch:b1:nucleophilic-substitution}

Ambillah sebotol \cip{S}-2-bromobutana, sebuah \omterm{def:b1:stereochemistry-in-depth:enantiomers}{enantiomer} tunggal yang
memutar cahaya terpolarisasi. Jika dipanaskan dengan natrium hidroksida
dalam \omterm{def:b1:intermolecular-forces-solvents:solvent-classes}{pelarut polar} aprotik, zat itu menghasilkan butan-2-ol yang juga
merupakan \omterm{def:b1:stereochemistry-in-depth:enantiomers}{enantiomer} tunggal --- tetapi \omterm{def:b1:stereochemistry-in-depth:enantiomers}{enantiomer} yang lain, \cip{R}:
reaksi itu telah membalik molekulnya seperti payung yang terbalik
ditiup angin. Jika dibiarkan dalam air hangat, bromida yang sama
menghasilkan butan-2-ol yang sebagian rasemat. Perubahan keseluruhan yang
sama, gugus OH menggantikan atom Br, mengikuti dua jalan yang berbeda.
Bab ini menetapkan kedua mekanisme itu dari \omterm{def:b1:rate-laws:order}{hukum lajunya}, menjelaskan
stereokimianya, dan memberikan aturan-aturan yang menentukan di antara
keduanya --- pemakaian penuh pertama alat-alat dari bab kinetika dan bab
efek elektronik.

\begin{recall}
\omterm{def:b1:rate-laws:order}{Hukum laju}, \omterm{def:b1:elementary-steps:elementary-step}{tahap elementer}, \omterm{def:b1:elementary-steps:rds}{tahap penentu laju}, dan \omterm{def:b1:elementary-steps:qssa}{pendekatan keadaan
tunak} terdapat pada \cref{ch:b1:rate-laws,ch:b1:elementary-steps};
\omterm{def:b1:stereochemistry-in-depth:isomers}{konfigurasi}, gambar Cram, dan \omterm{def:b1:stereochemistry-in-depth:optical-activity}{aktivitas optis} pada
\cref{ch:b1:stereochemistry-in-depth}; \omterm{def:b1:electronic-effects:nucleophile}{nukleofil}, \omterm{def:b1:electronic-effects:nucleophile}{elektrofil}, \omterm{def:b1:electronic-effects:nucleophile}{gugus pergi},
dan \omterm{def:b1:electronic-effects:intermediates}{karbokation} pada \cref{ch:b1:electronic-effects}; \omterm{def:b1:intermolecular-forces-solvents:solvent-classes}{pelarut protik} dan
aprotik pada \cref{ch:b1:intermolecular-forces-solvents}.
\end{recall}

\section{Substitusi dan para pelakunya}

\begin{definition}[Substitusi nukleofilik, substrat]\label{def:b1:nucleophilic-substitution:substitution}
\emph{Substitusi nukleofilik}\index{substitusi nukleofilik} adalah
penggantian, pada karbon jenuh, \omterm{def:b1:electronic-effects:nucleophile}{gugus pergi} Y oleh \omterm{def:b1:electronic-effects:nucleophile}{nukleofil} Nu:
$\mathrm{Nu^-} + \mathrm{R{-}Y} \to \mathrm{R{-}Nu} +
\mathrm{Y^-}$. Senyawa R--Y adalah \emph{substrat}\index{substrat};
karbon yang membawa Y bersifat elektrofilik karena ikatan C--Y
terpolarisasi ke arah Y.
\end{definition}

Haloalkana adalah \omterm{def:b1:nucleophilic-substitution:substitution}{substrat} yang khas, dengan \ce{Cl-}, \ce{Br-}, atau
\ce{I-} sebagai \omterm{def:b1:electronic-effects:nucleophile}{gugus pergi}. Hidroksida menghasilkan alkohol; alkoksida
menghasilkan eter; sianida menghasilkan nitril (bertambah satu karbon);
amonia dan amina menghasilkan amina; iodida bertukar dengan halida-halida
lain. Air dan alkohol, yang dipakai sebagai pelarut, juga dapat berperan
sebagai \omterm{def:b1:electronic-effects:nucleophile}{nukleofil}: substitusinya lalu disebut \omterm{def:b1:nucleophilic-substitution:sn1}{solvolisis}.

\section{Mekanisme SN2}

Ketika bromometana bereaksi dengan hidroksida di dalam air, menggandakan
salah satu konsentrasi menggandakan lajunya: $v = k[\ce{CH3Br}][\ce{OH-}]$,
reaksi berorde satu terhadap setiap pereaksi.

\begin{definition}[Mekanisme SN2, inversi Walden]\label{def:b1:nucleophilic-substitution:sn2}
\emph{Mekanisme SN2}\index{mekanisme SN2} (substitusi, nukleofilik,
bimolekuler) adalah satu \omterm{def:b1:elementary-steps:elementary-step}{tahap elementer} tunggal tempat \omterm{def:b1:electronic-effects:nucleophile}{nukleofil}
menyerang karbon dari sisi yang berlawanan dengan \omterm{def:b1:electronic-effects:nucleophile}{gugus pergi} sementara
gugus itu lepas. Ketiga substituen lain pada karbon itu terbalik,
seperti payung yang terbalik:
\emph{inversi Walden}\index{inversi Walden}.
\end{definition}

\begin{omfigure}
\setchemfig{atom sep=2.1em}
\schemestart
\chemfig{H@{nu}O^{-}}
\arrow{0}[,0.35]
\chemfig{@{c}C(-[:150]H_3C)(<[:210]H)(<:[:250]C_2H_5)-[@{cb}]@{br}Br}
\arrow{->}
\chemfig{[:0]HO-C(-[:30]CH_3)(<[:-30]H)(<:[:-70]C_2H_5)}
\+
\chemfig{Br^{-}}
\schemestop
\chemmove{\draw[->, thick, shorten <=2pt, shorten >=3pt](nu) .. controls +(15:5mm) and +(175:5mm) .. (c);
          \draw[->, thick, shorten <=2pt](cb) .. controls +(90:6mm) and +(90:6mm) .. (br);}

\medskip
\begin{tikzpicture}[font=\small]
  \node[draw, dashed, rounded corners, inner sep=6pt] at (0,0)
    {$\left[\vcenter{\hbox{\ce{HO}$\cdots$\chemfig{C(-[2]CH_3)(-[:-30]H)(-[:210]C_2H_5)}$\cdots$\ce{Br}}}\right]^{-}$};
  \node[below] at (0,-1.1) {keadaan transisi: lima gugus di sekitar C, tiga pada satu bidang};
\end{tikzpicture}

{\small Reaksi SN2 hidroksida dengan \cip{S}-2-bromobutana. \omterm{def:b1:electronic-effects:nucleophile}{Nukleofil}
menyerang sisi yang berlawanan dengan brom (\omterm{def:b1:electronic-effects:curly-arrow}{panah lengkung}); pada \omterm{def:b1:elementary-steps:energy-profile}{keadaan
transisi} ikatan O--C sedang terbentuk sementara ikatan C--Br putus;
produknya \cip{R}-butan-2-ol: ketiga gugus lain telah terbalik ke sisi
yang lain.}
\end{omfigure}

\begin{proposition}[Hukum laju SN2]\label{prop:b1:nucleophilic-substitution:sn2-law}
Untuk reaksi SN2, $v = k[\mathrm{RY}][\mathrm{Nu}]$.
\end{proposition}

\begin{proof}
Mekanismenya adalah satu \omterm{def:b1:elementary-steps:elementary-step}{tahap elementer} bimolekuler; menurut hukum aksi
massa untuk \omterm{def:b1:elementary-steps:elementary-step}{tahap elementer} (\cref{ch:b1:elementary-steps}) lajunya
sebanding dengan hasil kali konsentrasi kedua pihak.
\end{proof}

\begin{proposition}[Stereokimia SN2]\label{prop:b1:nucleophilic-substitution:sn2-inversion}
Reaksi SN2 pada karbon stereogenik menghasilkan satu \omterm{def:b1:stereochemistry-in-depth:isomers}{stereoisomer}
tunggal, dengan susunan ketiga gugus yang tidak berubah terbalik: reaksi
semacam ini bersifat stereospesifik.
\end{proposition}

\begin{proof}
\omterm{def:b1:electronic-effects:nucleophile}{Nukleofil} hanya dapat mencapai sisi karbon yang kosong, berlawanan
dengan \omterm{def:b1:electronic-effects:nucleophile}{gugus pergi}; ketiga gugus lain melewati sebuah bidang pada
\omterm{def:b1:elementary-steps:energy-profile}{keadaan transisi} dan berakhir di sisi yang lain. Ikatan baru mengarah ke
tempat yang tidak ditempati ikatan lama: setiap molekul \omterm{def:b1:nucleophilic-substitution:substitution}{substrat}
menghasilkan susunan gugus yang merupakan bayangan cerminnya.
Deskriptornya biasanya berubah dari \cip{S} menjadi \cip{R} atau
sebaliknya, tetapi tidak selalu, karena deskriptor itu bergantung pada
peringkat \omterm{def:b1:electronic-effects:nucleophile}{nukleofil} dan \omterm{def:b1:electronic-effects:nucleophile}{gugus pergi}.
\end{proof}

\begin{definition}[Reaksi stereospesifik]\label{def:b1:nucleophilic-substitution:stereospecific}
Sebuah reaksi disebut
\emph{reaksi stereospesifik}\index{reaksi stereospesifik} jika \omterm{def:b1:nucleophilic-substitution:substitution}{substrat-substrat} yang merupakan \omterm{def:b1:stereochemistry-in-depth:isomers}{stereoisomer}
menghasilkan produk-produk yang berbeda secara \omterm{def:b1:stereochemistry-in-depth:isomers}{stereoisomer}, setiap
\omterm{def:b1:nucleophilic-substitution:substitution}{substrat} menghasilkan produknya sendiri.
\end{definition}

\section{Mekanisme SN1}

2-Bromo-2-metilpropana bereaksi dengan air dengan laju yang tidak
bergantung pada konsentrasi \omterm{def:b1:electronic-effects:nucleophile}{nukleofil}: $v = k[(\ce{CH3})_3\ce{CBr}]$.
Sebuah tahap yang tidak melibatkan \omterm{def:b1:electronic-effects:nucleophile}{nukleofil} pasti menentukan lajunya.

\begin{definition}[Mekanisme SN1, rasemisasi, solvolisis]\label{def:b1:nucleophilic-substitution:sn1}
\emph{Mekanisme SN1}\index{mekanisme SN1} (substitusi, nukleofilik,
unimolekuler) mempunyai dua tahap: \omterm{def:b1:electronic-effects:cleavage}{heterolisis} ikatan C--Y yang lambat,
yang menghasilkan \omterm{def:b1:electronic-effects:intermediates}{karbokation}, lalu adisi \omterm{def:b1:electronic-effects:nucleophile}{nukleofil} yang cepat pada
\omterm{def:b1:electronic-effects:intermediates}{karbokation} itu. Reaksi yang mengubah satu \omterm{def:b1:stereochemistry-in-depth:enantiomers}{enantiomer} tunggal menjadi
campuran keduanya adalah \emph{rasemisasi}\index{rasemisasi};
substitusi yang \omterm{def:b1:electronic-effects:nucleophile}{nukleofilnya} adalah pelarut adalah
\emph{solvolisis}\index{solvolisis}.
\end{definition}

\begin{omfigure}
\setchemfig{atom sep=2.1em}
\schemestart
\chemfig{C(-[2]CH_3)(-[4]H_3C)(-[6]CH_3)-[@{cb2}]@{br2}Br}
\arrow{->[lambat]}
\chemfig{C^{+}(-[2]CH_3)(-[:210]H_3C)(-[:-30]CH_3)}
\+
\chemfig{Br^{-}}
\schemestop

\smallskip
\schemestart
\chemfig{C^{+}(-[2]CH_3)(-[:210]H_3C)(-[:-30]CH_3)}
\arrow{->[cepat][\ce{H2O}]}
\chemfig{C(-[2]CH_3)(-[4]H_3C)(-[6]CH_3)-OH_2^{+}}
\arrow{->[\ce{-H+}]}
\chemfig{C(-[2]CH_3)(-[4]H_3C)(-[6]CH_3)-OH}
\schemestop
\chemmove{\draw[->, thick, shorten <=2pt](cb2) .. controls +(90:6mm) and +(90:6mm) .. (br2);}

\medskip
\begin{tikzpicture}[font=\small]
  \fill[omProp!15, draw=omProp] (0,0.2) ellipse (0.2 and 0.85);
  \fill[omProp!15, draw=omProp] (0,-0.2) ellipse (0.2 and 0.85);
  \draw[thick] (0,0) -- (-1.1,-0.2) node[left] {\ce{CH3}};
  \draw[thick] (0,0) -- (0.95,-0.45) node[right] {\ce{C2H5}};
  \draw[thick] (0,0) -- (0.4,0.5) node[above right] {H};
  \node[fill=white, inner sep=1pt] at (0,0) {\ce{C+}};
  \draw[->, thick, omDef] (-0.9,1.6) node[left] {\ce{H2O}} -- (-0.15,0.95);
  \draw[->, thick, omDef] (0.9,-1.6) node[right] {\ce{H2O}} -- (0.15,-0.95);
  \node[align=center] at (4.6,0) {serangan pada kedua sisi:\\ kedua konfigurasi,\\ rasemat jika tidak ada lagi\\ yang membedakan kedua sisi};
\end{tikzpicture}

{\small \omterm{def:b1:nucleophilic-substitution:sn1}{Solvolisis} SN1 2-bromo-2-metilpropana di dalam air: ionisasi yang
lambat menjadi \omterm{def:b1:electronic-effects:intermediates}{karbokation} tersier, lalu penangkapan oleh air dan
pelepasan sebuah proton. Bawah: \omterm{def:b1:electronic-effects:intermediates}{karbokation} sekunder yang \omterm{def:b1:stereochemistry-in-depth:stereocentre}{kiral} (dari
2-bromobutana) berbentuk datar, dan air menyerang kedua sisinya dengan
peluang yang sama.}
\end{omfigure}

\begin{proposition}[Hukum laju SN1]\label{prop:b1:nucleophilic-substitution:sn1-law}
Untuk \omterm{def:b1:nucleophilic-substitution:sn1}{mekanisme SN1}, jika \omterm{def:b1:electronic-effects:intermediates}{karbokation} ditangkap jauh lebih cepat daripada
kembali menjadi \omterm{def:b1:nucleophilic-substitution:substitution}{substrat}, $v = k_1[\mathrm{RY}]$, tidak bergantung pada
\omterm{def:b1:electronic-effects:nucleophile}{nukleofil}.
\end{proposition}

\begin{proof}
Tahap-tahapnya: $\mathrm{RY} \rightleftharpoons \mathrm{R^+} + \mathrm{Y^-}$ ($k_1$,
$k_{-1}$), $\mathrm{R^+} + \mathrm{Nu} \to \mathrm{RNu}$ ($k_2$).
\omterm{def:b1:electronic-effects:intermediates}{Karbokation} adalah \omterm{def:b1:electronic-effects:intermediates}{zat antara reaktif}: menurut \omterm{def:b1:elementary-steps:qssa}{pendekatan keadaan tunak}
(\cref{ch:b1:elementary-steps}), $k_1[\mathrm{RY}] =
k_{-1}[\mathrm{R^+}][\mathrm{Y^-}] + k_2[\mathrm{R^+}][\mathrm{Nu}]$, sehingga
$v = k_2[\mathrm{R^+}][\mathrm{Nu}] = \dfrac{k_1k_2[\mathrm{RY}][\mathrm{Nu}]}
{k_{-1}[\mathrm{Y^-}] + k_2[\mathrm{Nu}]}$. Jika $k_2[\mathrm{Nu}] \gg
k_{-1}[\mathrm{Y^-}]$ ungkapan ini tereduksi menjadi $k_1[\mathrm{RY}]$:
tahap pertama adalah \omterm{def:b1:elementary-steps:rds}{tahap penentu laju}.
\end{proof}

\begin{proposition}[Stereokimia SN1]\label{prop:b1:nucleophilic-substitution:sn1-stereo}
Reaksi SN1 pada karbon stereogenik menghasilkan kedua \omterm{def:b1:stereochemistry-in-depth:isomers}{konfigurasi}:
reaksi itu tidak stereospesifik, dan menyebabkan \omterm{def:b1:nucleophilic-substitution:sn1}{rasemisasi} jika tidak
ada yang membedakan kedua sisi \omterm{def:b1:electronic-effects:intermediates}{karbokation}.
\end{proposition}

\begin{proof}
\omterm{def:b1:electronic-effects:intermediates}{Karbokation} itu datar (\cref{ch:b1:electronic-effects}): kedua sisinya
saling merupakan bayangan cermin. Jika \omterm{def:b1:electronic-effects:nucleophile}{gugus pergi} telah menjauh sebelum
penangkapan, \omterm{def:b1:electronic-effects:nucleophile}{nukleofil} menyerang kedua sisi dengan laju yang sama,
sehingga menghasilkan kedua \omterm{def:b1:stereochemistry-in-depth:enantiomers}{enantiomer} dalam jumlah yang sama. Dalam
praktik ion yang pergi masih menutupi satu sisi untuk beberapa saat, dan
sering teramati sedikit kelebihan inversi.
\end{proof}

\begin{omfigure}
\begin{tikzpicture}
\begin{axis}[
  width=0.8\linewidth, height=5.0cm, xmin=0, xmax=10, ymin=-30, ymax=110,
  axis x line=bottom, axis y line=left, xlabel={koordinat reaksi},
  ylabel={energi}, xtick=\empty, ytick=\empty,
  label style={font=\small}, legend style={font=\scriptsize, draw=none, at={(0.98,0.98)}, anchor=north east}]
\addplot[omDef, very thick] table[x=x, y=E] {figdata/out/bachelor-1-energy-profiles-sn2.dat};
\addplot[omProp, very thick, dashed] table[x=x, y=E] {figdata/out/bachelor-1-energy-profiles-sn1.dat};
\legend{SN2: satu tahap, SN1: dua tahap}
\node[font=\scriptsize, anchor=north] at (axis cs:5,68) {karbokation};
\node[font=\scriptsize, anchor=north] at (axis cs:0.6,-3) {RY};
\node[font=\scriptsize, anchor=south] at (axis cs:9.6,-18) {RNu};
\end{axis}
\end{tikzpicture}

{\small \omterm{def:b1:elementary-steps:energy-profile}{Profil energi} skematis. SN2: satu \omterm{def:b1:elementary-steps:energy-profile}{keadaan transisi}, dengan lima
gugus di sekitar karbon. SN1: sawar pertama yang lebih tinggi (ionisasi,
\omterm{def:b1:elementary-steps:rds}{tahap penentu laju}), \omterm{def:b1:electronic-effects:intermediates}{karbokation} dalam sebuah lembah yang dangkal, lalu
sawar kecil menuju penangkapannya.}
\end{omfigure}

\section{Penentu mekanisme}

\begin{proposition}[Substrat]\label{prop:b1:nucleophilic-substitution:substrate}
\omterm{def:b1:nucleophilic-substitution:substitution}{Substrat} metil dan primer bereaksi melalui SN2; \omterm{def:b1:nucleophilic-substitution:substitution}{substrat} tersier melalui
SN1; \omterm{def:b1:nucleophilic-substitution:substitution}{substrat} sekunder melalui salah satunya, bergantung pada
faktor-faktor lain. \omterm{def:b1:nucleophilic-substitution:substitution}{Substrat} alilik dan benzilik bereaksi cepat melalui
keduanya.
\end{proposition}

\begin{proof}
SN2 memerlukan ruang di belakang karbon: setiap gugus alkil yang
menggantikan hidrogen membuat \omterm{def:b1:elementary-steps:energy-profile}{keadaan transisi}, yang membawa lima gugus,
menjadi lebih padat dan memperlambat reaksinya; tiga gugus menghalanginya
sama sekali. SN1 memerlukan \omterm{def:b1:electronic-effects:intermediates}{karbokation} yang stabil: urutan kestabilan
(\cref{prop:b1:electronic-effects:carbocation-order}) membuat ionisasi
cepat untuk \omterm{def:b1:electronic-effects:intermediates}{karbokation} tersier, mungkin untuk sekunder, dan terlalu
lambat untuk \omterm{def:b1:electronic-effects:intermediates}{karbokation} primer dan metil. \omterm{def:b1:electronic-effects:intermediates}{Karbokation} alilik dan
benzilik distabilkan oleh resonansi, dan sistem $\pi$ yang sama
menstabilkan \omterm{def:b1:elementary-steps:energy-profile}{keadaan transisi} SN2.
\end{proof}

\begin{proposition}[Gugus pergi]\label{prop:b1:nucleophilic-substitution:leaving-group}
Kedua mekanisme makin cepat makin baik \omterm{def:b1:electronic-effects:nucleophile}{gugus perginya}, yaitu makin lemah
basa yang terbentuk darinya: $\ce{I-} > \ce{Br-} > \ce{Cl-} \gg \ce{F-}$;
\ce{OH-} dan \ce{RO-} tidak lepas.
\end{proposition}

\begin{proof}
Pada keduanya, ikatan C--Y putus pada atau sebelum \omterm{def:b1:elementary-steps:rds}{tahap penentu laju},
dan \omterm{def:b1:elementary-steps:energy-profile}{keadaan transisinya} membawa sebagian muatan negatif pada Y. Gugus
yang mengikat muatan negatif dengan baik --- basa konjugasi sebuah \omterm{def:b1:predominant-reaction:strong-acid}{asam
kuat} --- menurunkan \omterm{def:b1:elementary-steps:energy-profile}{keadaan transisi} itu. \ce{HI}, \ce{HBr}, dan
\ce{HCl} adalah \omterm{def:b1:predominant-reaction:strong-acid}{asam kuat}, \ce{HF} \omterm{def:b1:predominant-reaction:strong-acid}{asam lemah} ($\mathrm{p}K_a$ 3.2), air
dan alkohol asam yang sangat lemah: \ce{OH-} dan \ce{RO-} adalah \omterm{def:b1:predominant-reaction:strong-acid}{basa
kuat} dan \omterm{def:b1:electronic-effects:nucleophile}{gugus pergi} yang buruk.
\end{proof}

\begin{proposition}[Pelarut]\label{prop:b1:nucleophilic-substitution:solvent}
\omterm{def:b1:intermolecular-forces-solvents:solvent-classes}{Pelarut polar} protik (air, alkohol) mendukung SN1; \omterm{def:b1:intermolecular-forces-solvents:solvent-classes}{pelarut polar} aprotik
(propanon, dimetil sulfoksida, dimetilformamida) mendukung SN2 dengan
\omterm{def:b1:electronic-effects:nucleophile}{nukleofil} anionik.
\end{proposition}

\begin{proof}
SN1 menghasilkan dua ion dari \omterm{def:b1:nucleophilic-substitution:substitution}{substrat} netral: \omterm{def:b1:intermolecular-forces-solvents:solvent-classes}{pelarut polar} protik
menstabilkan keduanya, kation melalui \omterm{def:b1:lewis-resonance-vsepr:lewis}{pasangan elektron bebasnya} dan
anion melalui \omterm{def:b1:intermolecular-forces-solvents:hbond}{ikatan hidrogen}, serta menurunkan sawar ionisasi. Pada SN2
muatan \omterm{def:b1:electronic-effects:nucleophile}{nukleofil} anionik tersebar pada \omterm{def:b1:elementary-steps:energy-profile}{keadaan transisi}; \omterm{def:b1:intermolecular-forces-solvents:solvent-classes}{pelarut protik},
yang mengikat anion kecil itu dengan kuat melalui \omterm{def:b1:intermolecular-forces-solvents:hbond}{ikatan hidrogen},
memperlambatnya, sedangkan \omterm{def:b1:intermolecular-forces-solvents:solvent-classes}{pelarut aprotik} membiarkan anion itu bebas dan
reaktif (\cref{ch:b1:intermolecular-forces-solvents}).
\end{proof}

\begin{method}[Memilih mekanisme]\label{met:b1:nucleophilic-substitution:choose-mechanism}
\begin{enumerate}
  \item \omterm{def:b1:nucleophilic-substitution:substitution}{Substrat}: metil atau primer $\to$ SN2; tersier $\to$ SN1;
        sekunder $\to$ lanjutkan.
  \item \omterm{def:b1:electronic-effects:nucleophile}{Nukleofil}: kuat dan anionik ($\ce{OH-}$, $\ce{RO-}$, $\ce{CN-}$,
        $\ce{I-}$) $\to$ SN2; lemah dan netral (air, alkohol) $\to$ SN1.
  \item Pelarut: aprotik $\to$ SN2; protik $\to$ SN1.
  \item Ramalkan stereokimianya (inversi atau \omterm{def:b1:nucleophilic-substitution:sn1}{rasemisasi}) dan periksalah
        bahwa \omterm{def:b1:predominant-reaction:strong-acid}{basa kuat} tidak justru menyebabkan eliminasi
        (\cref{ch:b1:elimination}).
\end{enumerate}
\end{method}

\begin{omfigure}
\begin{tabular}{lll}
\hline
 & SN2 & SN1 \\
\hline
tahap & satu & dua, melalui \omterm{def:b1:electronic-effects:intermediates}{karbokation} \\
\omterm{def:b1:rate-laws:order}{hukum laju} & $k[\mathrm{RY}][\mathrm{Nu}]$ & $k[\mathrm{RY}]$ \\
\omterm{def:b1:nucleophilic-substitution:substitution}{substrat} & metil $>$ primer $>$ sekunder & tersier $>$ sekunder \\
\omterm{def:b1:electronic-effects:nucleophile}{nukleofil} & kuat, anionik & lemah, netral (sering kali pelarutnya) \\
pelarut & polar aprotik & polar protik \\
stereokimia & inversi (stereospesifik) & \omterm{def:b1:nucleophilic-substitution:sn1}{rasemisasi} (umumnya) \\
\hline
\end{tabular}

\smallskip
{\small Kedua mekanisme berdampingan.}
\end{omfigure}

\section{Latihan}

\begin{exercise}[$\star$]\label{exo:b1:nucleophilic-substitution:1}
Tuliskan produk-produk: bromoetana dengan natrium hidroksida; iodometana
dengan natrium etoksida; 1-kloropropana dengan kalium sianida; bromometana
dengan amonia (tahap pertama). Tentukan \omterm{def:b1:nucleophilic-substitution:substitution}{substrat}, \omterm{def:b1:electronic-effects:nucleophile}{nukleofil}, dan \omterm{def:b1:electronic-effects:nucleophile}{gugus
perginya}.
\end{exercise}

\begin{exercise}[$\star$]\label{exo:b1:nucleophilic-substitution:2}
Untuk reaksi $\mathrm{RY} + \mathrm{Nu}$ dengan $v = k[\mathrm{RY}][\mathrm{Nu}]$,
apa yang terjadi pada \omterm{def:b1:rate-laws:initial-rate}{laju awal} jika $[\mathrm{Nu}]$ dijadikan tiga kali?
Jika kedua konsentrasi dibuat setengahnya? Pertanyaan yang sama jika
$v = k[\mathrm{RY}]$.
\end{exercise}

\begin{exercise}[$\star$]\label{exo:b1:nucleophilic-substitution:3}
Ramalkan mekanismenya (SN1 atau SN2): 1-bromobutana dengan natrium iodida
dalam propanon; 2-bromo-2-metilpropana di dalam air; bromometana dengan
hidroksida; 2-bromobutana dalam etanol tanpa tambahan basa.
\end{exercise}

\begin{exercise}[$\star$]\label{exo:b1:nucleophilic-substitution:4}
Gambarlah produk reaksi SN2 sianida dengan \cip{R}-2-bromobutana dalam
\omterm{def:b1:stereochemistry-in-depth:representations}{representasi Cram} dan tentukan deskriptornya.
\end{exercise}

\begin{exercise}[$\star\star$]\label{exo:b1:nucleophilic-substitution:5}
Tuliskan \omterm{def:b1:nucleophilic-substitution:sn1}{mekanisme SN1} lengkap hidrolisis 2-bromo-2-metilpropana, dengan
\omterm{def:b1:electronic-effects:curly-arrow}{panah lengkung}, termasuk pelepasan protonnya. Mengapa \omterm{def:b1:rate-laws:order}{hukum lajunya}
berorde satu padahal tiga spesies terlibat?
\end{exercise}

\begin{exercise}[$\star\star$]\label{exo:b1:nucleophilic-substitution:6}
Jelaskan mengapa 1-bromo-2,2-dimetilpropana, sebuah bromida primer,
bereaksi sangat lambat melalui SN2 dan tidak bereaksi melalui SN1.
\end{exercise}

\begin{exercise}[$\star\star$]\label{exo:b1:nucleophilic-substitution:7}
\cip{S}-2-iodooktana kehilangan \omterm{def:b1:stereochemistry-in-depth:optical-activity}{aktivitas optisnya} jika disimpan dalam
propanon dengan natrium iodida, walaupun tidak terbentuk senyawa baru.
Jelaskan, lalu nyatakan mengapa laju hilangnya \omterm{def:b1:stereochemistry-in-depth:optical-activity}{aktivitas optis} dua kali
laju substitusinya.
\end{exercise}

\begin{exercise}[$\star\star$]\label{exo:b1:nucleophilic-substitution:8}
Urutkan menurut laju SN2 dengan iodida: 2-bromo-2-metilpropana,
bromometana, 2-bromopropana, 1-bromopropana. Urutkan senyawa-senyawa yang
sama menurut laju \omterm{def:b1:nucleophilic-substitution:sn1}{solvolisis} SN1.
\end{exercise}

\begin{exercise}[$\star\star$]\label{exo:b1:nucleophilic-substitution:9}
Sebuah bromida sekunder yang murni secara optis mengalami \omterm{def:b1:nucleophilic-substitution:sn1}{solvolisis};
produknya memutar cahaya sebesar \qty{12}{\percent} dari nilai yang
diharapkan untuk inversi murni. Tentukan perbandingan kedua \omterm{def:b1:stereochemistry-in-depth:enantiomers}{enantiomer}
produk itu dan tafsirkan.
\end{exercise}

\begin{exercise}[$\star\star\star$]\label{exo:b1:nucleophilic-substitution:10}
Turunkan \omterm{def:b1:rate-laws:order}{hukum laju} SN1 yang umum dengan \omterm{def:b1:elementary-steps:qssa}{pendekatan keadaan tunak}, lalu
tunjukkan bahwa menambahkan ion \omterm{def:b1:electronic-effects:nucleophile}{gugus pergi} \ce{Y-} pada awal reaksi
memperlambat reaksinya (\omterm{def:b1:precipitation:common-ion}{efek ion senama} dalam kinetika). Dengan syarat
apakah efek ini dapat diabaikan?
\end{exercise}

\begin{exercise}[$\star\star\star$]\label{exo:b1:nucleophilic-substitution:11}
Dua \omterm{def:b1:rate-laws:half-life}{waktu paruh} untuk reaksi bromometana dengan hidroksida: dengan
$[\ce{OH-}]_0 = [\ce{CH3Br}]_0 = C_0$, tunjukkan bahwa $1/[\ce{CH3Br}] = 1/C_0 +
kt$ dan tentukan $t_{1/2}$. Dengan $[\ce{OH-}]_0 = 50\,C_0$, tunjukkan
bahwa peluruhannya eksponensial dan tentukan $t_{1/2}$. Data latihan ini:
$k =
\qty{2.0e-4}{L/(mol.s)}$, $C_0 = \qty{0.010}{mol/L}$.
\end{exercise}

\begin{exercise}[$\star\star\star$]\label{exo:b1:nucleophilic-substitution:12}
Jelaskan mengapa sebuah alkohol, \ce{ROH}, tidak bereaksi dengan natrium
bromida, tetapi bereaksi dengan asam bromida pekat menghasilkan
\ce{RBr}. Tuliskan mekanismenya untuk alkohol tersier dan untuk alkohol
primer.
\end{exercise}

\section{Soal: Solvolisis tert-Butil Klorida}

\begin{problem}[{Soal akhir pekan --- orde dan tetapan laju dari data
konduktivitas, waktu paruh, mekanisme dan stereokimianya, serta energi
aktivasi solvolisis itu}]
\label{pb:b1:nucleophilic-substitution:1}
2-Kloro-2-metilpropana (tert-butil klorida) dilarutkan, pada konsentrasi
yang kecil, dalam campuran air--propanon. \omterm{def:b1:nucleophilic-substitution:sn1}{Solvolisisnya}
\ce{(CH3)3CCl + H2O -> (CH3)3COH + H+ + Cl-} menghasilkan ion, dan
konduktivitas larutan $\kappa$, yang sebanding dengan jumlah ion yang
terbentuk, direkam. Sesudah waktu yang sangat lama konduktivitas itu
mencapai $\kappa_\infty = \qty{2.000}{mS/cm}$ pada kedua suhu.
Hasil pengukuran (\unit{mS/cm}):

\begin{center}
\begin{tabular}{lcccccccc}
\hline
$t$ (min) & 0 & 5 & 10 & 15 & 20 & 30 & 45 & 60 \\
\hline
$\kappa$, \qty{25.0}{\celsius} & 0 & 0.172 & 0.329 & 0.473 & 0.605 & 0.835 & 1.110 & 1.321 \\
$\kappa$, \qty{35.0}{\celsius} & 0 & 0.507 & 0.885 & 1.168 & 1.379 & 1.654 & 1.856 & 1.940 \\
\hline
\end{tabular}
\end{center}
% data generated in figdata/bachelor-1/solvolysis.py (story data)

\medskip
\noindent\textbf{Bagian I --- Ordenya.}
\begin{enumerate}
  \item Mengapa konduktivitas sebanding dengan jumlah \omterm{def:b1:nucleophilic-substitution:substitution}{substrat} yang telah
        bereaksi?
  \item Nyatakan $[\mathrm{RCl}]/[\mathrm{RCl}]_0$ dalam $\kappa$ dan
        $\kappa_\infty$.
  \item Tuliskan hukum terintegrasi reaksi orde satu dan besaran yang
        harus digrafikkan terhadap $t$ untuk mengujinya.
  \item Hitunglah $\ln(\kappa_\infty - \kappa)$ pada \qty{25.0}{\celsius}
        untuk setiap waktu.
  \item Apakah grafiknya berupa garis lurus? Simpulkan ordenya.
  \item Air sangat berlebih. Apa yang tersembunyi dalam orde itu karena
        hal ini?
\end{enumerate}

\medskip
\noindent\textbf{Bagian II --- \omterm{def:b1:rate-laws:order}{Tetapan laju} dan \omterm{def:b1:rate-laws:half-life}{waktu paruh}.}
\begin{enumerate}[resume]
  \item Tentukan \omterm{def:b1:rate-laws:order}{tetapan laju} pada \qty{25.0}{\celsius}, dalam
        \unit{s^{-1}}.
  \item Hitunglah \omterm{def:b1:rate-laws:half-life}{waktu paruhnya}.
  \item Pada waktu berapakah \qty{90}{\percent} \omterm{def:b1:nucleophilic-substitution:substitution}{substrat} telah terpakai?
  \item Tentukan \omterm{def:b1:rate-laws:order}{tetapan laju} pada \qty{35.0}{\celsius}.
  \item Dengan faktor berapakah laju bertambah untuk sepuluh derajat?
  \item Periksalah satu nilai deret \qty{35}{\celsius} terhadap hukum
        itu.
\end{enumerate}

\medskip
\noindent\textbf{Bagian III --- Mekanismenya.}
\begin{enumerate}[resume]
  \item Mekanisme apakah yang diisyaratkan oleh \omterm{def:b1:nucleophilic-substitution:substitution}{substrat} dan pelarutnya?
  \item Tuliskan mekanisme itu dengan \omterm{def:b1:electronic-effects:curly-arrow}{panah lengkung}.
  \item Tunjukkan bahwa mekanisme itu sesuai dengan orde yang diperoleh.
  \item Menambahkan natrium hidroksida \qty{0.01}{mol/L} hampir tidak
        mengubah laju hilangnya klorida. Mengapa?
  \item Percobaan yang sama dengan klorida tersier yang \omterm{def:b1:stereochemistry-in-depth:stereocentre}{kiral},
        3-kloro-3-metilheksana, yang murni secara optis, menghasilkan
        alkohol dengan \omterm{def:b1:stereochemistry-in-depth:optical-activity}{aktivitas optis} hampir nol. Jelaskan.
  \item Gambarlah \omterm{def:b1:elementary-steps:energy-profile}{profil energi} reaksinya dan tandai \omterm{def:b1:elementary-steps:rds}{tahap penentu
        lajunya}.
  \item Mengapa 1-klorobutana hampir tidak bereaksi pada kondisi yang
        sama?
\end{enumerate}

\medskip
\noindent\textbf{Bagian IV --- \omterm{def:b1:rate-laws:activation-energy}{Energi aktivasi}.}
\begin{enumerate}[resume]
  \item Tuliskan \omterm{thm:b1:rate-laws:arrhenius}{hukum Arrhenius}.
  \item Nyatakan $E_a$ dari kedua \omterm{def:b1:rate-laws:order}{tetapan laju}.
  \item Hitunglah $E_a$ ($R = \qty{8.314}{J/(K.mol)}$).
  \item Apa yang diukur oleh energi ini, dalam mekanisme itu?
  \item Ramalkan \omterm{def:b1:rate-laws:order}{tetapan laju} pada \qty{45.0}{\celsius}.
  \item Nyatakan \omterm{def:b1:rate-laws:activation-energy}{energi aktivasi} \omterm{def:b1:nucleophilic-substitution:sn1}{solvolisis} itu, sampai dua angka
        penting.
\end{enumerate}
\end{problem}
